\subsection{An approximation bound and a discarded Dedekind construction}

\begin{lemma}[Proved: approximation by finitely generated subalgebras]
\label{A:A2-approximation}
Let $k$ be a field, let $\ell\ne\operatorname{char}k$, and suppose
$c=\operatorname{cd}_{\ell}k<\infty$.  Let $R$ be a $k$-algebra domain
whose fraction field $K$ has finite transcendence degree $t$ over $k$.
We do not assume that $R$ is noetherian or finitely generated.  Then
\[
\operatorname{cd}_{\ell}(\operatorname{Spec}R)\le c+t.
\]
In particular, if $k$ is algebraically closed, $K$ is algebraic over
$k(t)$, and $R$ is a Dedekind subring of $K$ with fraction field $K$,
then $R$ contains $k$, its $\ell$-cohomological dimension is at most
one, and $\operatorname{Pic}(R)/\ell=0$.
\end{lemma}

\begin{proof}
Write $R=\varinjlim_i R_i$, where $R_i$ runs through its finitely
generated $k$-subalgebras.  Every $R_i$ is a domain and
$\dim R_i=\operatorname{trdeg}_k\operatorname{Frac}(R_i)\le t$.
Artin's affine theorem over a separable closure of $k$, followed by
the Hochschild--Serre spectral sequence, gives
$\operatorname{cd}_{\ell}(\operatorname{Spec}R_i)\le c+t$; the
precise affine theorem is SGA~4, Expos\'e~XIV, Corollary~3.2
\cite{C_SGA4}.

Here are the two limit steps, including nonconstructible coefficients.
Every constructible $\mathbf Z/\ell^n$-sheaf on
$\operatorname{Spec}R$ descends to some $\operatorname{Spec}R_i$,
and its cohomology is the filtered colimit of the cohomology of its
pullbacks to later stages.  Consequently it has no cohomology in
degrees above $c+t$.  Every $\ell$-primary torsion sheaf is a filtered
colimit of constructible $\ell$-primary torsion sheaves, and cohomology
on the affine scheme $\operatorname{Spec}R$ commutes with that
filtered colimit.  We use, respectively, the Stacks Project,
Lemma~59.73.10(3), Theorem~59.51.3, Lemma~59.73.2(2)--(3), and
Lemma~59.51.4, all opened in this run \cite{C_StacksLimits}.
This proves the assertion for arbitrary torsion coefficients.

For the final assertion, consider a maximal ideal $\mathfrak m$ of
the Dedekind domain $R$ and its normalized discrete valuation
$v_{\mathfrak m}\colon K^\times\to\mathbf Z$.  The group $k^\times$
is divisible because $k$ is algebraically closed.  Its image in
$\mathbf Z$ is therefore zero.  Thus every nonzero element of $k$
is a unit in every $R_{\mathfrak m}$.  The intersection identity
$R=\bigcap_{\mathfrak m}R_{\mathfrak m}$ inside $K$ gives $k\subset R$.
If $R$ is itself a field, this inclusion is immediate.  We may now
apply the first assertion with $c=0$ and $t=1$.  Finally the Kummer
sequence injects $\operatorname{Pic}(R)/\ell$ into
$H^2(R,\mu_\ell)=0$.
\end{proof}

\begin{proposition}[Proved; the algebraic-number-field case of generic dimension two]
\label{A:A2-arithmetic-two}
Let $R$ be a characteristic-zero domain whose fraction field $K$ is
algebraic over $\Q$, and let $p$ be invertible in $R$. If $p=2$,
suppose that $K$ is not orderable. Then
\[
 \cd_p(\Spec R)\leq2.
\]
No noetherian hypothesis is needed for this upper bound. If $R$ is
noetherian and $\cd_p(K)=2$, then A.2 holds for $\Spec R$.
\end{proposition}
\begin{proof}
Write $R$ as the filtered union of its finitely generated
$\Z[1/p]$-subalgebras $R_i$. Each $F_i=\Frac(R_i)$ is a number
field. If $p=2$, choose one identity
$-1=\sum_j (a_j/b_j)^2$ in $K$, with $a_j,b_j\in R$ and
$b_j\ne0$, and restrict to the cofinal system containing all these
elements. Every resulting $F_i$ is then totally imaginary.

The finite-type one-dimensional domain $R_i$ has finite
normalization $\widetilde R_i$. Its spectrum is an open subscheme
of $\Spec\mathcal O_{F_i}$: a normal finite-type model of a number
field over $\Spec\Z$ is quasi-finite, and the birational form of
Zariski's main theorem identifies it with such an open. This open
omits the primes over $p$, because $p$ is a unit. Thus it is a
proper open of the arithmetic curve. Artin--Verdier duality in
\citet[Chapter II, Theorem 3.1, Corollary 3.2, and Proposition 2.3(a)]{A_Milne2006},
in the form verified in Proposition~\ref{A:arithmetic-test}, gives
$\cd_p(\widetilde R_i)\leq2$. For odd $p$ the real-place terms have
no positive $p$-cohomology; for $p=2$ there are no real places in
the chosen system.

Let $\nu:\Spec\widetilde R_i\to\Spec R_i$ be the finite
normalization, and let $\mathcal F$ be a constructible sheaf of
$\mathbf F_p$-vector spaces on $\Spec R_i$. The adjunction map
\[
 0\longrightarrow\mathcal F\longrightarrow
 \nu_*\nu^*\mathcal F\longrightarrow\mathcal Q\longrightarrow0
\]
is injective: on each geometric stalk it is the diagonal map into
the nonempty finite collection of stalks over that point. The
cokernel is supported on the finite closed set where $\nu$ is not
an isomorphism. Every residue field there is finite, so
$H^q(\Spec R_i,\mathcal Q)=0$ for $q>1$. Finite pushforward is
exact and preserves cohomology. The long exact sequence therefore
gives $H^q(\Spec R_i,\mathcal F)=0$ for $q>2$.
Coefficient d\'evissage gives the same bound for all $p$-primary
torsion sheaves.

Constructible descent, cohomological continuity, and approximation
of arbitrary torsion sheaves by constructible ones, with the exact
locators in Lemma~\ref{A:A2-approximation}, pass this uniform bound
to $R$. Every affine open of $\Spec R$ has the same type of
coordinate domain and the same fraction field. If its generic
dimension is two, Proposition~\ref{A:spread} supplies the matching
lower bound after shrinking, proving A.2.
\end{proof}

\begin{corollary}[Proved; why the completely split positive-$\mu$ towers do not refute A.2]
\label{A:A2-split-tower-exclusion}
Let $K_\infty/F$ be a $\Z_p$-extension of a number field, with $F$
totally imaginary if $p=2$. Suppose that some finite prime of $F$
not dividing $p$ splits completely in $K_\infty/F$. Then
$\cd_p(K_\infty)=2$. Every noetherian subring $R\subset K_\infty$
with fraction field $K_\infty$ and with $p$ invertible satisfies A.2.
The integral closure of $\mathcal O_F[1/p]$ in $K_\infty$ is not
noetherian.
\end{corollary}
\begin{proof}
Choose a compatible prime over the completely split prime.
The union of the corresponding discrete valuation rings is a
discrete valuation ring: all ramification indices and residue
degrees are one. Its residue field is the original finite field.
Lemma~\ref{A:valuation-lower} therefore gives
$\cd_p(K_\infty)\geq2$. The opposite inequality follows from
$\cd_p(F)=2$ and monotonicity under algebraic extension, with the
real-place convention as stated. Proposition~\ref{A:A2-arithmetic-two}
now applies to every such $R$ and to all its affine opens.

The integral closure has infinitely many distinct primes over the
chosen prime, since their number in the finite layers tends to
infinity. It is a one-dimensional domain integral over a Dedekind
domain. A nonzero element of the chosen prime of $F$ lies in all
these infinitely many maximal ideals. In a noetherian
one-dimensional domain the primes containing a given nonzero
element form a finite set, since its quotient is a zero-dimensional
noetherian ring. This is a contradiction.

The split-prime hypothesis is present in the constructions checked
in \citet[Proposition 3.1 and Corollary 3.2, pp.~11--12]{A_HajirMaire}
and \citet[Lemma 1 and \S3, Theorem 2]{A_HubbardWashington}.
The conclusion does not depend on their positive-$\mu$ calculation:
the split valuation alone forces generic dimension two, and
localizing the ring does not change that fraction field.
\end{proof}

\begin{lemma}[Proved; why a pseudo-algebraically closed fraction field cannot supply the proposed counterexample]
\label{A:A2-PAC-exclusion}
A pseudo-algebraically closed field admits no nontrivial discrete
valuation. Consequently it cannot be the fraction field of a
positive-dimensional noetherian domain.
\end{lemma}
\begin{proof}
Suppose $v:K^\times\to\Z$ is a surjective discrete valuation, and
choose $t\in K$ with $v(t)=1$. Choose a prime $q$ different from
$\operatorname{char}K$, and an integer $m\geq2$ also prime to that
characteristic. The affine curve
\[
 C:\quad y^q=t(1+t^q x^{mq})
\]
is geometrically integral. Indeed $1+t^q x^{mq}$ has simple nonzero
roots over $\overline K$. At any one of them the displayed monic
polynomial in $y$ is Eisenstein, hence irreducible over
$\overline K(x)$ and, by Gauss's lemma, over $\overline K[x,y]$.

It has no $K$-point. If $v(x)\geq0$, the valuation of its right-hand
side is one. If $v(x)<0$, it is
$1+q+mqv(x)$, again congruent to one modulo $q$. Neither is the
valuation of a nonzero $q$th power. The case $x=0$ belongs to the
first alternative, and the right-hand side is never zero for a
valued $K$-point. This contradicts the defining rational-point
property of a pseudo-algebraically closed field.

If a noetherian domain $R$ of positive dimension had fraction field
$K$, a height-one localization and then a localization of its
normalization would give a discrete valuation ring with fraction
field $K$. The normalization step uses the one-dimensional theorem
already opened at \citet[Tag 09IG, Lemma 10.120.18]{A_Stacks}; it
does not require finite normalization or excellence. This proves
the final assertion. The exclusion is classical; we include the
elementary curve argument to verify the failed construction route
without relying on an unopened textbook locator.
\end{proof}

Lemma~\ref{A:A2-approximation} rules out a persistent Picard-group obstruction obtained by
prescribing a Dedekind class group while keeping the fraction field
algebraic over a one-variable function field over an algebraically
closed field.  It does not rule out algebraic extensions of number
fields: there the cohomological dimension of the fraction field can
drop below the bound furnished by the original ground field.

\begin{corollary}[Proved under stated hypotheses; an additional A.2 case]
\label{A:A2-approximation-equality}
Let $R$ be a noetherian $k$-algebra domain, with $p\ne\operatorname{char}k$,
and suppose
\[
 \cd_p(\Frac R)=\cd_p(k)+\operatorname{trdeg}_k\Frac R<\infty.
\]
Then A.2 holds for $\Spec R$, without a finite-type or excellence
assumption on $R$.
\end{corollary}
\begin{proof}
Every affine open is the spectrum of a $k$-algebra domain with the
same fraction field. Lemma~\ref{A:A2-approximation} gives the upper
bound by the displayed common value on every such open.
Proposition~\ref{A:spread} gives the matching lower bound after
shrinking inside any prescribed nonempty open, and on every smaller
nonempty affine open.
\end{proof}

\begin{lemma}[Proved: finite character in an algebraic tower]
\label{A:A2-tower-Dedekind}
Let $R_0$ be a Dedekind domain and let $K_\infty/K_0$, with
$K_0=\operatorname{Frac}(R_0)$, be an algebraic extension expressed as
a directed union of finite extensions $K_i$.  Suppose that their
integral closures $R_i$ of $R_0$ are finite over $R_0$.  Suppose also
that each nonzero prime of every $R_i$ has only finitely many primes
above it in $R_\infty=\bigcup_iR_i$, and that along each such prime
the ramification indices in the tower are bounded.  Then
$R_\infty$ is a Dedekind domain.
\end{lemma}

\begin{proof}
Integrality shows that every nonzero prime of $R_\infty$ is maximal.
At such a prime $P$, the localization $(R_\infty)_P$ is the directed
union of the compatible valuation rings $(R_i)_{P\cap R_i}$.  It is
a valuation ring.  Bounded ramification indices make its value group
a discrete cyclic subgroup of $\mathbf Q$, so this localization is
a discrete valuation ring.  A nonzero element $a\in R_\infty$
belongs to some $R_i$ and has only finitely many prime divisors in
that $R_i$; the finite-splitting hypothesis then gives only finitely
many maximal ideals of $R_\infty$ containing $a$.

We check noetherianity directly.  Given a nonzero ideal $I$, choose
$0\ne a\in I$, and list the finitely many maximal ideals
$P_1,\dots,P_s$ containing $a$.  For each $j$ choose $b_j\in I$
whose valuation at $P_j$ is the minimum of the nonnegative integer
valuations attained on $I\setminus\{0\}$.  The ideal
$J=(a,b_1,\dots,b_s)$ has $J_{P_j}=I_{P_j}$ for every $j$; at every
other maximal ideal $a$ is a unit, so again $J_P=I_P$.  Local equality
of these ideals at all maximal ideals implies $I=J$.  Thus every
ideal is finitely generated.  The domain is noetherian, normal and
one-dimensional, as required.
\end{proof}
